\begin{abstract}
Tracking long-term capital accumulation and neighborhood-level economic change using survey-based datasets like the American Community Survey (ACS) is frequently constrained by high-variance sampling noise and transient macro-shocks. While aerial computer vision models offer high-resolution spatial tracking, they conventionally suffer from \textit{Structural Cardinality Mismatch}---memorizing the localized topological scale of their training data (e.g., New York's vertical density versus Dallas's horizontal sprawl) and failing to generalize out-of-sample due to temporal label drift. In this paper, we introduce an ordinal Siamese Network architecture explicitly designed with a dual objective: to map the cross-sectional distribution of structural wealth and to directionally rank longitudinal structural change (gentrification and decay) across a multi-year panel. 

To isolate true physical capital accumulation from transient econometric ``yo-yo'' shocks, we engineer a heavily smoothed, log-normalized longitudinal panel that strictly filters out sampling noise via temporal Z-tests. Rather than parametrically regressing cardinal income levels, our model acts as a pure non-parametric ordinal feature extractor. It is trained via a decoupled Triplet Margin Ranking Loss consisting of an unconditional L1 stability penalty for structurally unchanged buildings, and a signed directional hinge to rank temporal and spatial wealth shifts. By intentionally discarding topological cardinality and supervising only the direction of pairwise comparisons, the model prevents the label inconsistency inherent in longitudinal data. Ultimately, this purely ordinal framework enables zero-shot out-of-sample generalization to novel cities, tracking both static income distributions and dynamic urban evolution from visual infrastructure alone via direct empirical quantile matching with local survey moments.
\end{abstract}